\section{Reading and writing a first proof}\label{ch01:sec:first-proof}
Proofs use a few small words in a precise way. ``Let $a$ and $b$ be odd integers'' gives the names $a$ and $b$ to two odd integers that could be \emph{any} odd integers at all, not your two favorite examples. Since the proof may use only the fact that they are odd integers, the finished argument works whichever odd integers are supplied. Words such as ``hence'' and ``therefore'' announce a deduction, and the reader should always be able to see its reason.

Read the next result in three passes: first what is assumed and what is promised; then the plan, which shows how one might discover the proof but proves nothing itself; and then the proof, line by line.

\par\noindent\begin{minipage}{\linewidth}
\begin{proposition}\label{thm:odd-sum}
The sum of any two odd integers is even.
\end{proposition}
\end{minipage}\par

\noindent\textit{Plan.} Write each odd integer in the form given by the definition of odd, so that its oddness is visible. Then rearrange the sum until it has the form required by the definition of even: twice an integer.
\begin{proof}
Let $a$ and $b$ be odd integers. By the definition of odd (Definition~\ref{def:parity}), there are integers $k$ and $\ell$ with $a=2k+1$ and $b=2\ell+1$. The integers $k$ and $\ell$ may be equal or different; the argument does not depend on which. Substituting these expressions gives
\begin{align*}
a+b&=(2k+1)+(2\ell+1) &&\text{substitute for $a$ and $b$}\\
&=2k+2\ell+2 &&\text{reorder and regroup the terms, add $1+1$}\\
&=2(k+\ell+1) &&\text{factor out $2$.}
\end{align*}
The number $k+\ell+1$ is an integer, because a sum of integers is an integer. Hence $a+b$ is twice an integer. By the definition of even, $a+b$ is even.
\end{proof}

The final sentence connects the expression we computed to the property we promised: a proof ends not when the algebra looks tidy, but when the required conclusion has been reached.

To see what each symbol does, take $a=7$ and $b=-3$, with witnesses $k=3$ and $\ell=-2$. The proof builds the witness $k+\ell+1=2$ for the sum, and indeed $7+(-3)=4=2\cdot2$. This check is not a second proof.

The proposition promises an even sum only when \emph{both} numbers are odd; about $2+3=5$ it says nothing. Keep the hypotheses attached to a result when you remember it.

\begin{principle}{Habit 1: choose a representation that exposes the goal}
Before calculating, ask what form the answer must take. Then choose a way of writing the objects, a \emph{representation}, in which that form can appear. To show that a number is even, try to display a factor of $2$; to show that an approximation is close to the true value, write down the difference between them. When an assistant proposes a representation, ask why it helps before letting it continue.
\end{principle}
